\begin{statement}{C.2.UN.81}{Proved}{Finite removal of all vertical point defects at the marked Hopf interfaces}
Let \(f:X\to\Delta\) be a proper degeneration of connected
complex two-tori with smooth total space and full reduced
SNC central fibre. We also allow as the starting recorded model
the smooth-total NC endpoint obtained from such a fibre by
UN77, retaining its two labelled normalization branches at
each self-double curve. Normalize the actual canonical generator by
\[
 \operatorname{div}(\Omega)=\sum k_iV_i,\qquad
 \min k_i=0,\qquad M=\max k_i.
 \tag{C.AP.1}
\]
Suppose that the components of order \(M\) with Hopf minimal
model have already been made minimal by the actual record of
UN77. Mark any finite collection of their interfaces
\(J_\alpha=H_\alpha\cap P_\alpha\) for which the normal line of
\(J_\alpha\) is torsion and \(k_{P_\alpha}<M\).
Retain repeated markings on the same parent.

There is a finite overlap-compatible modification record,
unchanged over \(\Delta^*\), whose endpoint has smooth total
space and full reduced NC central fibre, such that every
designated parent is an embedded minimal fixed-Jacobian
elliptic quasi-bundle. Its full elliptic fibres are multiples
of smooth elliptic curves, its integral linear monodromy is
the identity, and every one of its boundary curves is a
smooth elliptic fibre support. In particular no triple point
meets a designated parent at the endpoint.

All maximal-order minimal Hopf components and an open
neighbourhood of their union are unchanged by the record.
In particular every marked \(J_\alpha\), its normal map and
its gluing to \(H_\alpha\) are retained. The designated parent
prime divisors survive. Other endpoint components can retain
NC self-incidences; their two normalization branches are
part of the record.

Every retained global component has canonical order between
\(0\) and the original \(M\). Whole-component removals delete
a component of positive order \(q\); their preliminary
additions, when needed, have order \(q+1\le M\), and there
are no additions at \(q=M\). The record retains its actual
centres, local maps, normalization branches and restored
overlaps, so reversing it recovers the original family.
For a given compact family the same modifications are
performed on that compact base. Their Euler changes are
the exact normal-cone terms of UN75.

No assertion is made here that every maximal Hopf component
has one of the marked interfaces. In particular an
equal-order interface with a nontorsion normal line is not
among these markings.
\end{statement}

\begin{proof}
\par\medskip\noindent\textbf{The exact marked parent maps.}
All curves on a minimal Hopf surface are disjoint smooth
elliptic curves. Thus each marked curve is isolated from
the other boundaries on both sides. Its actual normal-product
isomorphism makes the normal on the parent torsion as well.
On the normalization of a designated parent \(P\), of order
\(p\), the marked curve has coefficient
\(M-p-1\ge0\) in the actual canonical divisor. UN70 therefore
gives a specified point map and elliptic map
\[
 P\xrightarrow{\pi_P} Z_P\xrightarrow{q_P}B_P,
 \qquad R^1(q_P)_*\OO_{Z_P}=\OO_{B_P},
 \qquad g(B_P)\le1.
 \tag{C.AP.2}
\]
The point map is an isomorphism on a neighbourhood of every
marked curve. Every boundary is vertical for \(q_P\pi_P\),
and every rational boundary is exceptional for \(\pi_P\).
The exact constant-period field of UN80, equations
(C.EP.1)--(C.EP.11), proves the additional inclusion
\[
                 \operatorname{Exc}(\pi_P)\subset D_P.
 \tag{C.AP.3}
\]
Here \(D_P\) is the full actual reduced boundary. It includes
both preimages of a self-double curve at an NC stage.

These assertions also apply at the intermediate smooth-total
NC stages used below. Indeed the actual adjunction divisor
on the normalization counts those two preimages; the marked
elliptic neighbourhood stays embedded and unchanged. The
saturated ambient fields restrict as meromorphic fields to
that normalization and are holomorphic off its full boundary.
The surface maps in UN70 and the direct exceptional-valuation
argument in UN80 consequently apply to this actual smooth
normalization, without identifying its two boundary branches.

The target \(Z_P\) is a minimal surface in its birational
class. Thus
\[
 h_P=b_2(P)-b_2(Z_P)\ge0
 \tag{C.AP.4}
\]
is exactly the number of point blowups in \(\pi_P\). The
minimal second Betti number is unchanged under the subsequent
birational maps: in the ruled case it is two, and in the
other cases the minimal model is unique. We count each
designated parent once in \(h=\sum_P h_P\).

\par\medskip\noindent\textbf{The full fibre graph and the receiving degrees.}
The point map \(\pi_P\) is a sequence of point blowups of
smooth elliptic fibre supports. For one such fibre, retain
its entire reduced inverse image, whether or not its strict
elliptic component is an actual boundary. Its graph is an
elliptic root with rational trees. Write \(C^2=-b_C\) and let
\(n_C\) be the number of its intersections with the other
components of this entire reduced fibre. The point maps give
\[
 n_C\le b_C+1\quad\hbox{for a rational component},\qquad
 n_C\le b_C\quad\hbox{for the elliptic component}.
 \tag{C.AP.5}
\]
The elliptic root starts with \((b,n)=(0,0)\). A new rational
curve has \((b,n)=(1,1)\) at a smooth point and \((1,2)\) at
a node. Blowing up a smooth point of an old component raises
both \(b\) and \(n\) by one; a node blowup raises \(b\) by one
and leaves its valence unchanged. These explicit updates
prove (C.AP.5) at every infinitely-near point.

Choose a last point blowup of \(\pi_P\), with exceptional
curve \(E\). By (C.AP.3), \(E\) is an actual boundary curve.
It meets zero, one or two other actual boundary curves.
These are exactly its triple points in the threefold.
The complete fibre graph has at most two neighbours at this
last point, so there is no additional case.

Suppose first that there are two. On the other component
\(Q\), the normal-product formula gives \(E_Q^2=-1\), because
\(E_P^2=-1\). Hence the actual type-II two-ruling diagram of
UN75 applies. It blows down \(E\) on both adjacent surfaces
and blows up the two opposite normalized branches at their
specified points.

Neither receiving branch belongs to a designated parent.
To prove this, let \(C\) be the opposite boundary on \(P\),
and write \(-b_0\) and \(n_0\) for its square and entire-fibre
valence just before the last node blowup. Afterwards its
square is \(-b_0-1\), and its number \(t_C\) of actual triple
points is at most \(n_0\). On the opposite receiving surface
\(R\), the exact normal map therefore gives
\[
 C_R^2=b_0+1-t_C\ge b_0+1-n_0\ge0,
 \quad\hbox{and } C_R^2\ge1\hbox{ if }C\hbox{ is elliptic}.
 \tag{C.AP.6}
\]
On a designated parent every rational boundary is exceptional
and has negative square; every elliptic boundary is the
strict transform of a fibre support and has nonpositive
square. Both alternatives contradict (C.AP.6). This argument
uses the individual normalized branch when the two receivers
are branches of one global component. Thus a type-II move
leaves the component counts unchanged and strictly decreases
\(h\); it does not create a self-double curve on a designated
parent.

\par\medskip\noindent\textbf{One triple point: an actual ruled-removal block.}
Suppose that \(E\) has one triple point. Then \(E_Q^2=0\).
Its complete rational deformations give the actual ruling
on the compact normalization of \(Q\), after its vertical
point contractions. The boundary passing through this one
triple point is its unique section; every other boundary is
vertical. This is the same complete-curve and actual ruling
comparison used in (C.HM.9), not a ruling inferred from a
numerical class alone. If \(q=k_Q\) and \(R\) is the section
component, restriction of the actual canonical divisor to
a complete ruling fibre gives
\[
                         k_R=q-1,
                         1\le q\le M.
 \tag{C.AP.7}
\]
The actual finite ruled-removal and cleanup block in UN77
therefore applies. On a projective-line bundle its contraction
is precisely UN74; for nonminimal fibres it is the retained
sigma-atlas construction of
\cite[\S\S4.10.7--4.10.9,4.14]{Kulikov1977}, with the
missing-fibre correction of \cite[\S8]{Kulikov1981}.

For clarity, the numerical and analytic properties of this
block used here are the ones already proved in UN75 and
UN77. It removes a positive-order component \(q\), adds
finitely many components of order \(q+1\) when necessary,
and adds none when \(q=M\). The diagnostic resolution used
to check Lemma4.14 is not inserted into the component list.
Its actual sigma-atlas is cleaned up by choosing a largest
last local branch, following its finite descendants and
performing the same ruled removals. Empty lists give smooth
NC charts. All overlaps are restored by the restrictions of
the same proper local graph maps. The positive component
card vector strictly decreases across this whole block.
No monotonicity of \(h\) inside the block is required.

\par\medskip\noindent\textbf{Zero triple points: the exact plane comparison.}
Suppose that \(E\) has no triple point. Its square on \(Q\)
is \(+1\), and every other boundary on \(Q\) is disjoint
from \(E\). Here rationality does not require \(Q\) to be
a canonical maximum. A rational curve of positive square
cannot be exceptional for a minimal-model point map. On a
minimal model of nonnegative Kodaira dimension, the explicit
canonical lines in UN31 and its quotient classification are
nef. Their pullback and the positive point discrepancies
would give \(K_Q\cdot E\ge0\), contrary to adjunction
\(K_Q\cdot E=-3\). A Hopf minimal model contains no rational
curve. On an irrational ruled minimal model every rational
curve is vertical for its ruling and has nonpositive square,
and point blowups cannot increase that square. The component
classification therefore leaves exactly the rational case.
One can specify its required plane
map without choosing a birational presentation. The sequence
\[
 0\longrightarrow\OO_Q\longrightarrow\OO_Q(E)
       \longrightarrow\OO_E(1)\longrightarrow0
 \tag{C.AP.8}
\]
and \(H^1(Q,\OO_Q)=0\) give three sections. They have no base
point on \(E\), and the section defining \(E\) is nonzero
off \(E\). Thus \(|E|\) gives an actual morphism
\[
 \rho:Q\longrightarrow\PP^2,
 \qquad \rho^*\OO_{\PP^2}(1)=\OO_Q(E),
 \qquad \rho|_E:E\xrightarrow{\sim}\ell .
 \tag{C.AP.9}
\]
Its image is a surface since \(E^2=1\); its degree is one
by the same intersection number. Hence it is a birational
morphism, factoring into the actual point blowdowns away
from \(E\). Every other boundary is contracted, since it
has zero intersection with \(E\).

The canonical restriction to \(E\) gives
\(k_Q=k_P+2\). More generally call a current pair \((Q,E)\)
of the form (C.AP.9), with all other boundary disjoint from
\(E\), a marked plane pair, and write \(q=k_Q\). Its actual
canonical divisor has the precise form
\[
 K_Q=-3E+\sum_C a_CC,
 \qquad \sum_C a_CC=K_Q-\rho^*K_{\PP^2},
 \qquad a_C>0.
 \tag{C.AP.10}
\]
Indeed choose on \(\PP^2\) the meromorphic two-form whose
divisor is \(-3\ell\). The ratio of the actual form on \(Q\)
to its pullback has divisor supported on the exceptional
curves. On \(\PP^2\) this same meromorphic function has no
divisorial zero or pole, so it is a nonzero constant. Thus
the equality in (C.AP.10) is an equality of actual divisors,
with the positive discrepancies of the specified point map.
It shows that every exceptional prime is an actual
boundary, and that there is no additional boundary of
coefficient zero. The only coefficient on the marked curve
is \(-3\). A self-double curve away from that marking would
have coefficient \(-1\), so such a pair stays embedded at
the smooth-total NC stages under consideration.

Choose a last point blowup of \(\rho\). Its exceptional
boundary curve \(L\) lies away from \(E\). At a smooth point
of a discrepancy component of coefficient \(a>0\) its new
coefficient is \(a+1\); at a node of coefficients \(a,b>0\)
it is \(a+b+1\); at a point off the discrepancy divisor it
is one. These are the actual point-map identities for
\(K_Q+3E\). They give the following exhaustive alternatives:
\[
\begin{array}{c|c|c}
 \text{triple points of }L&\text{order of its neighbour}
                       &\text{actual next operation}\\ \hline
 0&q+2&\text{another marked plane pair}\\
 1&q+a+2\ge q+3&\text{a ruled-removal block}\\
 2&q+a+b+2&\text{a type-II move}.
\end{array}
 \tag{C.AP.11}
\]
In the last row the two receiving orders are
\(q+a+1\) and \(q+b+1\), both at least \(q+2\). The type-II
move decreases \(b_2(Q)\) by one; every other increase in
surface second Betti number occurs at a strictly higher
order. Its centre and all its points on \(Q\) are away from
the marked \(E\), so (C.AP.9) persists on the modified \(Q\).
In the first row the same complete linear system constructs
the next marked plane pair, at a strictly greater order.

This yields a finite progress block. Until its first whole
component removal, let
\[
                    b_j=\sum_{k_V=j}b_2(V^\nu),
                    \qquad 0\le j\le M,
 \tag{C.AP.12}
\]
where the sum uses the actual smooth normalizations. Each
type-II step in (C.AP.11) strictly decreases
\((b_0,\ldots,b_M)\) lexicographically. Passing to another
marked plane pair makes no modification and strictly raises
its order. There can be at most \(M\) such successive calls
before an actual step. Since lexicographic order on this
fixed finite product of nonnegative integers is well founded,
these steps cannot continue indefinitely without a component
removal. At a marked plane pair with no exceptional curve,
the surface is \(\PP^2\), its sole boundary is \(E=\ell\),
and the principal full-fibre restriction gives
\[
                 N_{Q/X}=\OO_{\PP^2}(-1).
 \tag{C.AP.13}
\]
The actual point contraction of UN64 removes it to a smooth
point. The adjacent surface is its ordinary point blowup,
and no other component passes through that point. In the
one-triple alternative the positive ruled-removal block
already supplies a component removal.

Stop this progress block at its first whole-component
removal, and perform the finite sigma cleanup just described.
Its net positive card vector decreases. It is unnecessary
to claim that the first marked plane pair or its marking
survives an arbitrary later cleanup: after this completed
block the outer argument uses the actual designated parent
maps again. This distinction retains every map while avoiding
an unsupported global preservation assertion.

\par\medskip\noindent\textbf{Preserved maximal Hopf neighbourhoods.}
Every type-II centre used on a designated parent is vertical
for (C.AP.2), hence disjoint from its marked elliptic curves.
An opposite curve of a type-II move cannot lie on a minimal
maximal Hopf: a rational curve is impossible there, and the
elliptic alternative has the positive square (C.AP.6).
For the plane progress moves all opposite boundary curves
are rational. They likewise cannot meet a minimal Hopf.

A ruled component removed in one of the actual blocks has
one boundary section, whose neighbouring order is \(q-1\),
and all other boundary is rational and vertical. It cannot
meet a minimal Hopf of order \(M\): that Hopf contains no
rational curve, while its being the section parent would
require \(q=M+1\). The same statement applies to every
cleanup removal and missing-fibre insertion. Plane point
contractions also have rational boundary. Thus these actual
closed centres and local sigma charts avoid the entire
union of maximal minimal Hopfs. Their restored overlap maps
are the identity on an open neighbourhood of that union.

The designated parent prime divisors cannot be removed:
each meets one of these unchanged Hopfs. Their marked
elliptic neighbourhoods therefore persist. At every smooth
NC stage the exact maps (C.AP.2)--(C.AP.4) can be applied
again to their current normalizations. No holomorphic
contraction of a Hopf neighbourhood has been invoked.

\par\medskip\noindent\textbf{The fixed finite termination potential.}
Let \(c_j\) count the retained global components of order
\(j\), for \(1\le j\le M\). Use the block-level potential
\[
                  (c_1,c_2,\ldots,c_M,h).
 \tag{C.AP.14}
\]
All coordinates are nonnegative integers. A direct type-II
step on a last designated-parent point leaves the cards
unchanged and decreases \(h\) by (C.AP.6). Either of the
other last-point alternatives supplies a finite progress
block with a strict decrease of the first changed card.
Any finite defect increase inside that whole block is
irrelevant to this lexicographic decrease. The missing-fibre
and cleanup comparisons keep all retained orders at most
the original \(M\); their auxiliary common blowups are not
additional retained cards. Hence (C.AP.14) has no infinite
strictly descending sequence.

When this procedure stops, a positive designated-parent
defect would still have a last point curve. UN80 and the
three exhaustive alternatives would supply another strict
decrease. Thus every such defect is zero. Each designated
normalization is the actual minimal elliptic model in
(C.AP.2), and every boundary is a fibre support. All new
self-incidences created by the type-I/type-II and ruled
sigma charts are rational exceptional curves. None can
remain on these parents, since their rational boundary
would have to be exceptional for a point map of defect
zero. No elliptic self-incidence was introduced: a ruled
removal has only one horizontal boundary section, and all
its other changes are at its rational fibre curves. The
designated parents are consequently embedded, with disjoint
elliptic boundaries and no triple point.

Every arrow used above is one of the actual two-ruling,
curve, point, or restored sigma-atlas maps already specified.
The inverse record, including all normalization branches
and overlap restrictions, reconstructs the given family.
The proper graph gluing and annular comparison of UN75
give the same record on a compact base, and its normal-cone
Euler formula applies at every retained step. This proves
the stated marked-interface reduction in its exact scope.
\end{proof}


% C.2.UN.57 from part_C2_all_diagonal_cycles.tex

